\documentclass[11pt]{article}
\usepackage[a4paper,margin=18mm]{geometry}
\usepackage[no-math]{fontspec}
\setmainfont{Latin Modern Roman}
\usepackage{amsmath,amssymb,amsthm,xfrac,xspace,hyperref,graphicx,comment}
\excludecomment{DefTralics}
\providecommand{\backmatter}{}

\providecommand{\bibliofont}{\small}
\newcommand{\ie}{i.e.,\xspace}
\newcommand{\eg}{e.g.,\xspace}
\newcommand{\etc}{etc.\xspace}
\newcommand{\cf}{cf.\xspace}
\newcommand{\resp}{resp.\xspace}
\newcommand{\smaller}{\small}
\newcommand{\Subsection}{\subsection}
\newcommand{\Subsubsection}{\subsubsection}
\newcommand{\skpt}{\smallskip}
\emergencystretch=3em
\numberwithin{equation}{section}\input{author-preamble.tex}
\begin{document}
\begin{center}{\Large Stable item 2056}\end{center}
\paragraph{Source and attribution.}
Yves Benoist, \emph{Positive harmonic functions on the Heisenberg group II}, Journal de l’École polytechnique — Mathématiques 8 (2021), publisher version, DOI 10.5802/jep.163. \href{https://jep.centre-mersenne.org/articles/10.5802/jep.163/}{Publisher article}. Authors' text: \href{https://creativecommons.org/licenses/by/4.0/}{CC BY 4.0}.
\paragraph{Editorial problem boundary.}
Prove only Theorem 1.1 below for the discrete Heisenberg group and a positive measure with finite support generating the group. The support need not generate a semigroup and the measure need not be normalized. Every extremal positive harmonic function is proportional to a character or a right translate of a function induced from a harmonic character of an abelian subgroup, exactly as defined in the source. All new induced/noninduced recognition, negligible-trajectory, switching and central-invariance arguments are retained. The finer existence/count-of-induced-functions classification and nilpotent rank-four counterexample remain surrounding context only.
\paragraph{Difficulty (editorial): H2.}
The cited Choquet--Deny theorem only classifies the abelian quotient or subgroup once central invariance or induction has been established. For general group-generating finite support, the author must prove that each extremal Heisenberg harmonic function falls into one of those two classes, including the case where the projected cone contains no line. The new weighted switching, negligible-trajectory and partition-generating-function estimates supply that nonroutine dichotomy before the final abelian reduction.
\paragraph{Editorial transcription note.}
Original author language, mathematics and ordering are preserved. Publisher front matter is replaced by this independent article wrapper. Bibliography is recovered from matching cached publisher HTML, including its TeX formulas. No new proof or mathematical repair is supplied. Surrounding original results are retained for conservative context and are not extra selected corpus claims.
\clearpage
\input{author-body.tex}
\begin{thebibliography}{99}
\bibitem{Ancona90} A. Ancona - “Théorie du potentiel sur les graphes et les variétés” , in École d’été de Probabilités de Saint-Flour XVIII—1988 , Lect. Notes in Math. , vol. 1427 , Springer, Berlin, 1990, p. 1-112
\bibitem{BenoistHeisenberg1} Y. Benoist - “Positive harmonic functions on the Heisenberg group I” , 2019
\bibitem{Breuillard05} E. Breuillard - “Local limit theorems and equidistribution of random walks on the Heisenberg group” , Geom. Funct. Anal. 15 (2005) no. 1, p. 35-82
\bibitem{Breuillard10} E. Breuillard - “Equidistribution of dense subgroups on nilpotent Lie groups” , Ergodic Theory Dynam. Systems 30 (2010) no. 1, p. 131-150
\bibitem{ChoquetDeny60} G. Choquet \& J. Deny - “Sur l’équation de convolution $\mu =\mu \ast \sigma $” , C. R. Acad. Sci. Paris 250 (1960), p. 799-801
\bibitem{DiaconisHough15} P. Diaconis \& B. Hough - “Random walk on unipotent matrix groups” , 2015
\bibitem{DynkinMaljutov61} E. B. Dynkin \& M. B. Maljutov - “Random walk on groups with a finite number of generators” , Dokl. Akad. Nauk SSSR 137 (1961), p. 1042-1045
\bibitem{GollSchmidtVerbitskiy16} M. Göll, K. Schmidt \& E. Verbitskiy - “A Wiener lemma for the discrete Heisenberg group” , Monatsh. Math. 180 (2016) no. 3, p. 485-525
\bibitem{Gouezel15} S. Gouëzel - “Martin boundary of random walks with unbounded jumps in hyperbolic groups” , Ann. Probab. 43 (2015) no. 5, p. 2374-2404
\bibitem{Guivarch71} Y. Guivarc’h - “Groupes nilpotents et probabilité” , C. R. Acad. Sci. Paris Sér. A-B 273 (1971), p. A997-A998
\bibitem{LindSchmidt15} D. Lind \& K. Shmidt - “A survey of algebraic actions of the discrete Heisenberg group” , Uspekhi Mat. Nauk 70 (2015) no. 4(424), p. 77-142
\bibitem{Margulis66} G. A. Margulis - “Positive harmonic functions on nilpotent groups” , Soviet Math. Dokl. 7 (1966), p. 241-244
\bibitem{Sawyer97} S. A. Sawyer - “Martin boundaries and random walks” , in Harmonic functions on trees and buildings (New York, 1995) , Contemp. Math. , vol. 206 , American Mathematical Society, Providence, RI, 1997, p. 17-44
\bibitem{VershikMalyutin18} A. M. Vershik \& A. V. Malyutin - “Asymptotics of the number of geodesics in the discrete Heisenberg group” , Zap. Nauchn. Sem. S.-Peterburg. Otdel. Mat. Inst. Steklov. (POMI) 468 (2018) no. XXIX, p. 39-52
\bibitem{Woess00} W. Woess - Random walks on infinite graphs and groups , Cambridge Tracts in Math. , vol. 138 , Cambridge University Press, Cambridge, 2000
\end{thebibliography}
\end{document}
